% GENERATED FILE. Edit canonical lessons, then run npm run generate:books.
% canonical-source-hash: fnv1a64:eb9d91047a3d98d5
% canonical-lessons: JA-C388-ryohi, JA-C388-koutsuuhi, JA-C388-tsuuka, JA-C388-muryou, JA-C388-yuuryou

\chapter{What the Journey Costs}
\label{ch:ja-what-the-journey-costs}
\hlchaptermodality{388}

\begin{chapteropening}
\textbf{By the end of this chapter:} \emph{I can talk about what a journey costs and what is free.}

Complete the last lesson of chapter 388: I can talk about what a journey costs and what is free.
\end{chapteropening}

\section[\texorpdfstring{\ja{りょひ} (ryohi)}{りょひ (ryohi)}]{\ja{りょひ} (ryohi) — travel expenses}

\label{lesson:JA-C388-ryohi}



\begin{quote}
Before the new one: say the Japanese for a plane ticket, then the Japanese for a business trip.
\end{quote}

\subsection*{You'll want to know: \ja{りょひ}}
\textbf{\ja{りょひ}} — \emph{ryohi} — \textquotedblleft{}travel expenses\textquotedblright{}.

\begin{quote}
\textbf{\ja{りょこうの} \ja{りょひは} \ja{ちちが} \ja{はらいました。}} — \emph{ryokō no ryohi wa chichi ga haraimashita} — \textquotedblleft{}My father paid the travel expenses for the trip.\textquotedblright{}
\end{quote}

\begin{cousinweb}
\textbf{\ja{りょひ}} (\emph{ryohi}) is Sino-Japanese (kango): two Chinese characters read with their on'yomi, readings borrowed from Chinese. They are \emph{ryo}, \textquotedblleft{}trip\textquotedblright{} (Middle Chinese \emph{ljoX}) and \emph{hi}, \textquotedblleft{}expense\textquotedblright{} (\emph{phj+jH}) (KANJIDIC2; Middle Chinese in Baxter's transcription).
\end{cousinweb}

\begin{grammarlens}[title={What you've built}]
One more word for getting around.
\end{grammarlens}

\subsection*{Your turn}
\begin{itemize}
\raggedright
  \item \emph{Say it:} \emph{ryohi}
  \item \emph{Say it:} \emph{ryohi}, once more
  \item \emph{Say it:} \emph{ryokō no ryohi wa chichi ga haraimashita} — My father paid the travel expenses for the trip.
  \item \emph{Recall:} say the Japanese for a plane ticket, then the Japanese for a business trip, then say \emph{ryohi} again
\end{itemize}

\begin{tcolorbox}[breakable,colback=teal!4,colframe=teal!35!black,title={Before you move on}]
What does \ja{りょひ} mean? (\textquotedblleft{}Travel expenses\textquotedblright{}.) Say it once more.
\end{tcolorbox}
\section[\texorpdfstring{\ja{こうつうひ} (kōtsūhi)}{こうつうひ (kōtsūhi)}]{\ja{こうつうひ} (kōtsūhi) — transport costs}

\label{lesson:JA-C388-koutsuuhi}



\begin{quote}
Before the new one: say the Japanese for a business trip, then the Japanese for travel expenses.
\end{quote}

\subsection*{You'll want to know: \ja{こうつうひ}}
\textbf{\ja{こうつうひ}} — \emph{kōtsūhi} — \textquotedblleft{}transport costs\textquotedblright{}.

\begin{quote}
\textbf{\ja{かいしゃが} \ja{こうつうひを} \ja{はらいます。}} — \emph{kaisha ga kōtsūhi o haraimasu} — \textquotedblleft{}The company pays the transport costs.\textquotedblright{}
\end{quote}

\begin{cousinweb}
\textbf{\ja{こうつうひ}} (\emph{kōtsūhi}) is Sino-Japanese (kango): three Chinese characters read with their on'yomi, readings borrowed from Chinese. They are \emph{kō}, \textquotedblleft{}coming \& going\textquotedblright{} (Middle Chinese \emph{kaew}), \emph{tsū}, \textquotedblleft{}traffic\textquotedblright{} (\emph{thuwng}) and \emph{hi}, \textquotedblleft{}expense\textquotedblright{} (\emph{phj+jH}) (KANJIDIC2; Middle Chinese in Baxter's transcription).
\end{cousinweb}

\begin{grammarlens}[title={What you've built}]
One more word for getting around.
\end{grammarlens}

\subsection*{Your turn}
\begin{itemize}
\raggedright
  \item \emph{Say it:} \emph{kōtsūhi}
  \item \emph{Say it:} \emph{kōtsūhi}, once more
  \item \emph{Say it:} \emph{kaisha ga kōtsūhi o haraimasu} — The company pays the transport costs.
  \item \emph{Recall:} say the Japanese for a business trip, then the Japanese for travel expenses, then say \emph{kōtsūhi} again
\end{itemize}

\begin{tcolorbox}[breakable,colback=teal!4,colframe=teal!35!black,title={Before you move on}]
What does \ja{こうつうひ} mean? (\textquotedblleft{}Transport costs\textquotedblright{}.) Say it once more.
\end{tcolorbox}
\section[\texorpdfstring{\ja{つうか} (tsūka)}{つうか (tsūka)}]{\ja{つうか} (tsūka) — a currency}

\label{lesson:JA-C388-tsuuka}



\begin{quote}
Before the new one: say the Japanese for travel expenses, then the Japanese for transport costs.
\end{quote}

\subsection*{You'll want to know: \ja{つうか}}
\textbf{\ja{つうか}} — \emph{tsūka} — \textquotedblleft{}a currency\textquotedblright{}.

\begin{quote}
\textbf{\ja{がいこくの} \ja{つうかを} \ja{すこし} \ja{もっています。}} — \emph{gaikoku no tsūka o sukoshi motte imasu} — \textquotedblleft{}I have a little foreign currency.\textquotedblright{}
\end{quote}

\begin{cousinweb}
\textbf{\ja{つうか}} (\emph{tsūka}) is Sino-Japanese (kango): two Chinese characters read with their on'yomi, readings borrowed from Chinese. They are \emph{tsū}, \textquotedblleft{}pass through\textquotedblright{} (Middle Chinese \emph{thuwng}) and \emph{ka}, \textquotedblleft{}goods\textquotedblright{} (\emph{xwaH}) (KANJIDIC2; Middle Chinese in Baxter's transcription).
\end{cousinweb}

\begin{grammarlens}[title={What you've built}]
One more word for getting around.
\end{grammarlens}

\subsection*{Your turn}
\begin{itemize}
\raggedright
  \item \emph{Say it:} \emph{tsūka}
  \item \emph{Say it:} \emph{tsūka}, once more
  \item \emph{Say it:} \emph{gaikoku no tsūka o sukoshi motte imasu} — I have a little foreign currency.
  \item \emph{Recall:} say the Japanese for travel expenses, then the Japanese for transport costs, then say \emph{tsūka} again
\end{itemize}

\begin{tcolorbox}[breakable,colback=teal!4,colframe=teal!35!black,title={Before you move on}]
What does \ja{つうか} mean? (\textquotedblleft{}A currency\textquotedblright{}.) Say it once more.
\end{tcolorbox}
\section[\texorpdfstring{\ja{むりょう} (muryō)}{むりょう (muryō)}]{\ja{むりょう} (muryō) — free of charge}

\label{lesson:JA-C388-muryou}



\begin{quote}
Before the new one: say the Japanese for transport costs, then the Japanese for a currency.
\end{quote}

\subsection*{You'll want to know: \ja{むりょう}}
\textbf{\ja{むりょう}} — \emph{muryō} — \textquotedblleft{}free of charge\textquotedblright{}.

\begin{quote}
\textbf{\ja{こどもの} \ja{うんちんは} \ja{むりょうです。}} — \emph{kodomo no unchin wa muryō desu} — \textquotedblleft{}Children travel free.\textquotedblright{}
\end{quote}

\begin{cousinweb}
\textbf{\ja{むりょう}} (\emph{muryō}) is Sino-Japanese (kango): two Chinese characters read with their on'yomi, readings borrowed from Chinese. They are \emph{mu}, \textquotedblleft{}nothing\textquotedblright{} (Middle Chinese \emph{mju}) and \emph{ryō}, \textquotedblleft{}fee\textquotedblright{} (\emph{lew}) (KANJIDIC2; Middle Chinese in Baxter's transcription).
\end{cousinweb}

\begin{grammarlens}[title={What you've built}]
One more word for getting around.
\end{grammarlens}

\subsection*{Your turn}
\begin{itemize}
\raggedright
  \item \emph{Say it:} \emph{muryō}
  \item \emph{Say it:} \emph{muryō}, once more
  \item \emph{Say it:} \emph{kodomo no unchin wa muryō desu} — Children travel free.
  \item \emph{Recall:} say the Japanese for transport costs, then the Japanese for a currency, then say \emph{muryō} again
\end{itemize}

\begin{tcolorbox}[breakable,colback=teal!4,colframe=teal!35!black,title={Before you move on}]
What does \ja{むりょう} mean? (\textquotedblleft{}Free of charge\textquotedblright{}.) Say it once more.
\end{tcolorbox}
\section[\texorpdfstring{\ja{ゆうりょう} (yūryō)}{ゆうりょう (yūryō)}]{\ja{ゆうりょう} (yūryō) — requiring a fee; not free}

\label{lesson:JA-C388-yuuryou}



\begin{quote}
Before the new one: say the Japanese for a currency, then the Japanese for free of charge.
\end{quote}

\subsection*{You'll want to know: \ja{ゆうりょう}}
\textbf{\ja{ゆうりょう}} — \emph{yūryō} — \textquotedblleft{}requiring a fee; not free\textquotedblright{}.

\begin{quote}
\textbf{\ja{えきの} \ja{ちゅうしゃじょうは} \ja{ゆうりょうです。}} — \emph{eki no chūshajō wa yūryō desu} — \textquotedblleft{}The station car park charges a fee.\textquotedblright{}
\end{quote}

\begin{cousinweb}
\textbf{\ja{ゆうりょう}} (\emph{yūryō}) is Sino-Japanese (kango): two Chinese characters read with their on'yomi, readings borrowed from Chinese. They are \emph{yū}, \textquotedblleft{}have\textquotedblright{} (Middle Chinese \emph{hjuwX}) and \emph{ryō}, \textquotedblleft{}fee\textquotedblright{} (\emph{lew}) (KANJIDIC2; Middle Chinese in Baxter's transcription).
\end{cousinweb}

\begin{grammarlens}[title={What you've built}]
One more word for getting around.
\end{grammarlens}

\subsection*{Your turn}
\begin{itemize}
\raggedright
  \item \emph{Say it:} \emph{yūryō}
  \item \emph{Say it:} \emph{yūryō}, once more
  \item \emph{Say it:} \emph{eki no chūshajō wa yūryō desu} — The station car park charges a fee.
  \item \emph{Recall:} say the Japanese for a currency, then the Japanese for free of charge, then say \emph{yūryō} again
\end{itemize}

\begin{tcolorbox}[breakable,colback=teal!4,colframe=teal!35!black,title={Before you move on}]
What does \ja{ゆうりょう} mean? (\textquotedblleft{}Requiring a fee; not free\textquotedblright{}.) Say it once more.
\end{tcolorbox}
